\documentclass[11pt,a4paper]{article}
\usepackage[utf8]{inputenc}
\usepackage[T1]{fontenc}
\usepackage{amsmath,amsfonts,amssymb}
\usepackage{graphicx}
\usepackage{hyperref}
\usepackage{listings}
\usepackage{color}
\usepackage{booktabs}
\usepackage{float}
\usepackage{geometry}
\geometry{margin=1in}
\definecolor{zenblue}{RGB}{41,121,255}
\hypersetup{colorlinks=true,linkcolor=zenblue,urlcolor=zenblue,citecolor=zenblue}

\title{\textbf{Zen3-Nano: A Compact-Deployment Finetune of Qwen3-8B}\\
\large Technical Whitepaper v2025.06}
\author{Zach Kelling \\ Zen LM Research Team\\
\texttt{research@zenlm.org}\\
\href{https://papers.zenlm.org}{papers.zenlm.org}}
\date{June 2025}

\begin{document}
\maketitle

\begin{abstract}
Zen3-Nano is a packaging- and deployment-focused member of the Zen family. Despite the ``Nano''
name, the released Zen3-Nano weights are \emph{not} a 0.6B-parameter from-scratch model: configuration
fingerprinting shows an 8B-class dense transformer matching \textbf{Qwen3-8B} \cite{qwen3report}, the
open-weight model released by Alibaba's Qwen team under the Apache-2.0 license. Earlier versions of
this whitepaper described a bespoke ``Zen MoDE (Mixture of Distilled Experts)'' architecture, a 0.6B
parameter count, a from-scratch 2-trillion-token pretraining run, and a hierarchical-distillation
pipeline from ``Zen3-Mini/Small/Medium'' teachers. Those claims do not match the shipped artifact and
have been removed. This corrected whitepaper documents the actual Qwen3-8B base, the light
instruction tuning and quantization-for-deployment work the Zen team performs, and defers to the
upstream Qwen3 technical report for pretraining and benchmark numbers. Quantized GGUF/MLX builds are
provided for convenient deployment; the genuine ``small'' tiers of the Qwen3 family (0.6B, 1.7B, 4B)
are noted as the appropriate bases for true edge use.
\end{abstract}

\tableofcontents
\newpage

%% -----------------------------------------------------------------------
\section{Introduction}
\label{sec:intro}

On-device and resource-constrained inference is an important deployment target, and quantization
makes even 8B-class models broadly deployable. This whitepaper documents Zen3-Nano honestly,
correcting prior inaccurate claims about its size and origin.

\paragraph{Provenance and corrections.} The released Zen3-Nano weights fingerprint to
\textbf{Qwen3-8B}: 36 layers, hidden size 4096, 32 query heads / 8 KV heads (GQA), head dimension
128, FFN intermediate size 12{,}288, and a 151{,}936-token vocabulary \cite{qwen3hf}. This is an
8.2B-parameter dense model released by the Qwen team at Alibaba Cloud under the Apache-2.0 license
\cite{qwen3report}. The following earlier claims were inaccurate and are withdrawn:

\begin{itemize}
\item ``0.6B parameters'' / ``sub-250MB'' / ``third-generation 0.6B edge model'' --- the shipped
      model is 8B-class, not 0.6B.
\item ``Zen MoDE (Mixture of Distilled Experts)'' architecture --- Qwen3-8B is a standard dense
      transformer, not a mixture-of-experts model; there is no ``MoDE'' routing.
\item ``Trained from scratch on a 2T-token corpus'' and ``hierarchical distillation from
      Zen3-Mini/Small/Medium teachers'' --- no such from-scratch pretraining or teacher cascade was
      performed by the Zen team; the base is the released Qwen3-8B.
\item Specific benchmark figures (e.g., MMLU 47.3\%, GSM8K 52.4\%, and the per-hardware
      tokens/second table) --- these were fabricated and are removed.
\end{itemize}

\paragraph{What Zen3-Nano actually is.} Zen3-Nano is a lightly instruction-tuned and quantized
repackaging of Qwen3-8B aimed at easy local deployment. Its contributions are (i) optional light SFT
for formatting/instruction style on top of Qwen3-8B, and (ii) ready-to-run quantized artifacts
(GGUF for llama.cpp, MLX for Apple Silicon, ONNX). For genuinely tiny edge footprints, the Qwen3
family also offers true small dense models---Qwen3-0.6B, Qwen3-1.7B, and Qwen3-4B (all Apache-2.0)
\cite{qwen3report}---which are the appropriate bases when sub-1B size is a hard requirement.

%% -----------------------------------------------------------------------
\section{Architecture (Inherited from Qwen3-8B)}
\label{sec:arch}

Zen3-Nano inherits the Qwen3-8B architecture unchanged: a dense decoder-only transformer with
grouped-query attention, SwiGLU feed-forward blocks, RMSNorm, and rotary position embeddings. There
is no mixture-of-experts component.

\begin{table}[H]
\centering
\caption{Architecture hyperparameters, inherited unchanged from Qwen3-8B \cite{qwen3hf}.}
\begin{tabular}{lc}
\toprule
\textbf{Hyperparameter} & \textbf{Value} \\
\midrule
Parameters (total)         & 8.2B \\
Non-embedding parameters   & 6.95B \\
Layers                     & 36 \\
Attention heads (query)    & 32 \\
KV heads (GQA)             & 8 \\
Head dimension             & 128 \\
Hidden dimension           & 4096 \\
FFN intermediate dimension & 12{,}288 \\
Vocabulary size            & 151{,}936 \\
Context length (native)    & 40{,}960 \\
Context length (with YaRN) & 131{,}072 \\
Position encoding          & RoPE ($\theta = 1{,}000{,}000$) \\
Activation function        & SiLU (SwiGLU) \\
Normalization              & RMSNorm \\
Tied embeddings            & No \\
\bottomrule
\end{tabular}
\label{tab:arch}
\end{table}

\subsection{Grouped-Query Attention}

Qwen3-8B uses GQA \cite{ainslie2023gqa} with 32 query heads grouped onto 8 key-value heads (head
dimension 128), reducing KV-cache memory relative to full multi-head attention. Rotary positional
embeddings (RoPE) \cite{su2021rope} are used with base frequency $\theta = 1{,}000{,}000$; long
context up to 131{,}072 tokens is available via YaRN scaling \cite{peng2023yarn} as configured
upstream.

\subsection{Feed-Forward Network}

Each layer uses a SwiGLU feed-forward block \cite{shazeer2020glu}:

\begin{equation}
\text{FFN}(\mathbf{x}) = \left(\text{SiLU}(\mathbf{x}\mathbf{W}_{\text{gate}}) \odot
\mathbf{x}\mathbf{W}_{\text{up}}\right) \mathbf{W}_{\text{down}}
\end{equation}

with intermediate dimension 12{,}288. There is no expert routing; the previously described
``$\text{MoDE}(\mathbf{x}) = \sum_{i} g_i(\mathbf{x}) E_i(\mathbf{x})$'' formulation does not apply
to this model.

%% -----------------------------------------------------------------------
\section{Post-Training and Quantization}
\label{sec:training}

The Zen team does not pretrain Zen3-Nano. Pretraining of the underlying Qwen3-8B was performed by the
Qwen team; see the upstream technical report \cite{qwen3report} for the corpus, scale, and procedure.
Our work on top of the released weights is limited to:

\subsection{Light Instruction Tuning (Optional)}

Where a chat-tuned variant is shipped, we apply light supervised fine-tuning on curated
instruction--response data (response-masked loss, low learning rate, few epochs) to adjust formatting
and instruction-following style. This does not constitute pretraining and does not materially change
the base model's knowledge.

\subsection{Quantization for Deployment}

The deployment artifacts use standard post-training quantization (e.g., llama.cpp's Q4\_K\_M and
Q8\_0 schemes, and MLX 4-bit) applied to the released weights. We do not perform the
``quantization-aware training from epoch 1'' described in earlier versions of this whitepaper; the
shipped artifact is the released Qwen3-8B quantized post hoc with widely used GGUF/MLX recipes.
References to a proprietary ``BitDelta (ZIP-007)'' quantization-aware objective have been removed
as they were not used to produce this model.

%% -----------------------------------------------------------------------
\section{Evaluation}
\label{sec:eval}

\paragraph{On benchmark numbers.} This whitepaper no longer reports fabricated benchmark tables or
per-device throughput figures. Because Zen3-Nano is the Qwen3-8B model (optionally lightly tuned and
quantized), its capabilities are those of Qwen3-8B, subject to the usual accuracy cost of 4-bit
quantization. For rigorous, independently reproducible results on Qwen3 models (MMLU, GSM8K, MATH,
HumanEval, multilingual, and long-context), see the official Qwen3 technical report
\cite{qwen3report} and the Qwen3-8B model card \cite{qwen3hf}. Any quantization-vs-BF16 quality
deltas we publish are measured on the released GGUF/MLX artifacts under a stated harness; we do not
quote numbers we have not measured.

\paragraph{Right-sizing for edge.} If a true sub-1B footprint is required, the appropriate choice is
a genuinely small Qwen3 base (0.6B/1.7B/4B), not an 8B model relabeled ``Nano.'' Published Qwen3
results for those tiers \cite{qwen3report} should be used to set expectations rather than the
fabricated figures previously printed here.

%% -----------------------------------------------------------------------
\section{Related Work}
\label{sec:related}

Knowledge distillation \cite{hinton2015distilling} and teacher-assistant cascades
\cite{mirzadeh2020improved} are well-established techniques for producing smaller models; however,
Zen3-Nano as shipped is not the product of such a pipeline---it is the Qwen3-8B base, optionally
lightly tuned and quantized. Post-training quantization for efficient deployment is standard practice
(e.g., the GGUF schemes used by llama.cpp). The underlying base model belongs to the Qwen3 series
\cite{qwen3report}, which provides a coherent family of dense (0.6B--32B) and MoE models under
Apache-2.0.

%% -----------------------------------------------------------------------
\section{Deployment and Runtime}
\label{sec:deploy}

Because Zen3-Nano is architecturally Qwen3-8B, it runs on any Qwen3-compatible runtime:

\begin{itemize}
\item \textbf{llama.cpp} --- GGUF format (Q4\_K\_M, Q8\_0).
\item \textbf{MLX} --- Apple Silicon, 4-bit.
\item \textbf{ONNX Runtime} --- ARM and x86 SIMD.
\item \textbf{Hugging Face Transformers / vLLM} --- BF16 reference weights.
\end{itemize}

\subsection{Context Length}

Native context is 40{,}960 tokens (Qwen3-8B default), extensible to 131{,}072 tokens via YaRN at the
cost of additional KV-cache memory. Memory footprint depends on quantization: an 8B model in 4-bit
requires on the order of several gigabytes, which is larger than the sub-250MB figure incorrectly
cited in earlier versions.

\subsection{API Compatibility}

When served via a Qwen3-compatible inference server, Zen3-Nano exposes an OpenAI-compatible API,
enabling drop-in replacement for OpenAI SDK clients.

%% -----------------------------------------------------------------------
\section{Limitations}
\label{sec:limitations}

Zen3-Nano inherits the limitations of Qwen3-8B and of autoregressive models generally: factual
hallucination, degraded performance on long multi-step reasoning, and incomplete safety alignment.
Post-training quantization to 4-bit introduces an additional, measurable accuracy cost relative to
BF16. The ``Nano'' name is a legacy product label and does not indicate a sub-1B model; deployments
with hard memory limits should select a genuinely small Qwen3 base instead.

%% -----------------------------------------------------------------------
\section{Conclusion}
\label{sec:conclusion}

Zen3-Nano, as released, is the Apache-2.0 Qwen3-8B model---optionally lightly instruction-tuned and
quantized for convenient local deployment---not a 0.6B from-scratch ``Zen MoDE'' model. This
corrected whitepaper documents the inherited Qwen3-8B architecture, the limited post-training and
standard quantization actually performed, and the correct guidance to use a true small Qwen3 base
(0.6B/1.7B/4B) when a sub-1B edge footprint is required. We defer to the upstream Qwen3 technical
report for pretraining details and rigorous, attributed benchmark numbers.

\section*{Acknowledgments}

We thank the Qwen team at Alibaba Cloud for releasing the Qwen3 models under the Apache-2.0 license,
and the maintainers of llama.cpp and MLX for the quantization and runtime tooling used to produce the
deployment artifacts.

\bibliographystyle{plain}
\begin{thebibliography}{9}

\bibitem{qwen3report}
Qwen Team, Alibaba Cloud,
``Qwen3 Technical Report,''
\textit{arXiv:2505.09388}, 2025.

\bibitem{qwen3hf}
Qwen Team,
``Qwen3-8B model card and configuration,''
Hugging Face, \url{https://huggingface.co/Qwen/Qwen3-8B}, 2025.

\bibitem{hinton2015distilling}
G. Hinton, O. Vinyals, and J. Dean,
``Distilling the Knowledge in a Neural Network,''
\textit{NIPS Deep Learning Workshop}, 2015.

\bibitem{mirzadeh2020improved}
S. I. Mirzadeh et al.,
``Improved Knowledge Distillation via Teacher Assistant,''
\textit{AAAI}, 2020.

\bibitem{ainslie2023gqa}
J. Ainslie et al.,
``GQA: Training Generalized Multi-Query Transformer Models,''
\textit{EMNLP}, 2023.

\bibitem{su2021rope}
J. Su et al.,
``RoFormer: Enhanced Transformer with Rotary Position Embedding,''
\textit{arXiv:2104.09864}, 2021.

\bibitem{shazeer2020glu}
N. Shazeer,
``GLU Variants Improve Transformer,''
\textit{arXiv:2002.05202}, 2020.

\bibitem{peng2023yarn}
B. Peng et al.,
``YaRN: Efficient Context Window Extension of Large Language Models,''
\textit{arXiv:2309.00071}, 2023.

\end{thebibliography}

\end{document}
